% ==============================================================
\section{はじめに}\label{sec:intro}
% ==============================================================

頻出パターンマイニング~\cite{agrawal1994apriori,han2000mining}は
トランザクションデータ中に繰り返し出現するアイテムセットを発見する基盤技術であり、
小売分析、医療データ解析、Web ログ分析など幅広い分野に応用されている。

Apriori-Window~\cite{aprioriwindow}は、
タイムスタンプ付きトランザクションデータベースに対して、
スライディングウィンドウのサポートが閾値 $\theta$ 以上となる
\emph{密集区間}とそこで頻出するパターンを効率的に発見するアルゴリズムである。
同手法は $t$ 軸ストライド最適化と Apriori 枝刈りを組み合わせることで、
候補区間の評価コストを大幅に削減し、
LPFIM や LPPM などの既存手法を性能で上回ることを示している。

しかし現実のトランザクションデータは時刻だけでなく\textbf{空間的位置情報}
（店舗座標、地域コード、GPS 位置等）を伴うことが多い。
例えば、小売チェーンでは各トランザクションが特定の店舗（座標）で発生し、
疫学データでは症例の発生場所が記録される。
このような\emph{時空間トランザクションデータ}において、
「沿岸部店舗のみで集中的に売れるアイテムセット」のような
\textbf{時空間局所的な密集パターン}を発見することは、
時間軸のみを扱う Apriori-Window では不可能である。
全店舗データを時間軸に集約してしまうと、
空間的に局在するパターンが背景ノイズに埋もれてしまうからである。

\subsection{技術的課題}

Apriori-Window を1次元（時間軸）から3次元（時間 $\times$ 空間 $x \times$ 空間 $y$）に
拡張するには以下の課題がある。

\begin{enumerate}
\item \textbf{支持曲面の計算爆発}：
グリッド $(T \times X \times Y)$ 上の全パターン・全グリッド点でサポートを
評価する単純実装は $O(M \cdot G \cdot W_t W_x W_y)$ の計算量を要する。
$G = T \times X \times Y$ はグリッドサイズ、
$M$ はパターン数、$W = W_t W_x W_y$ はウィンドウ体積である。

\item \textbf{次元増加による候補爆発}：
空間次元が加わると候補グリッド点が $T \to T \times X \times Y$ に増大し、
各パターンについてすべての点を評価するコストが膨大になる。

\item \textbf{$t$ 軸ストライドの3次元化}：
Apriori-Window の1次元ストライドアルゴリズム~\cite{aprioriwindow}を
各空間ウィンドウ $(x_s, y_s)$ に独立に適用できるよう構造化する必要がある。

\item \textbf{密集領域定義の一般化}：
1次元の「密集区間」を3次元の「密集領域（連結成分）」に自然に拡張し、
Apriori 枝刈りを密集セル集合の交差として定式化する必要がある。
\end{enumerate}

\subsection{本研究の貢献}

本稿では Apriori-Window~\cite{aprioriwindow}を
時間 $\times$ 空間 $(x, y)$ の3次元グリッドに拡張した
\textbf{3次元 Apriori-Window アルゴリズム}を提案する（\ref{sec:method}節）。
貢献は以下の3点である。

\begin{enumerate}
\item \textbf{3次元接頭辞和による $O(1)$ ウィンドウクエリ}：
出現グリッド $I_P$ に対して $(T{+}1) \times (X{+}1) \times (Y{+}1)$ 接頭辞和を構築し、
8コーナー包除原理により任意ウィンドウのサポートを $O(1)$ で取得する。
これは1次元の累積和クエリを3次元に拡張したものである。

\item \textbf{$t$ 軸ストライドの3次元拡張}（$k = 1$）：
Apriori-Window の $t$ 軸ストライドアルゴリズム~\cite{aprioriwindow}を
各空間ウィンドウ $(x_s, y_s)$ の XY イベントリストに独立に適用することで、
不要な時間スキャンを排除する（補題~\ref{lem:stride}で正確性を証明）。

\item \textbf{3次元 Apriori 枝刈り}（$k \geq 2$）：
Apriori-Window の候補区間交差枝刈り~\cite{aprioriwindow}を
3次元密集セル集合の交差として拡張する（補題~\ref{lem:apriori}）。
交差集合が空のパターンは $I_P$ 構築を省略でき、大幅な枝刈りが可能である。
\end{enumerate}

合成データ実験（\ref{sec:experiments}節）では、
3次元設定（$W_x = W_y = 3$）が植え付け3パターンを Precision $= 1.0$、
Recall $= 1.0$ で完全回収する一方、
1次元相当の設定（$W_x = W_y = 50$）では背景ノイズ由来の偽陽性パターンが
多数検出される（Precision $\approx 0.02$）ことを示し、
空間次元の明示的なモデル化の重要性を実証する。
また $N = 100{,}000$ トランザクションで約 1.5 秒（Rust + Rayon 並列化）を達成し、
Apriori 枝刈りが $k \geq 3$ の候補を完全に排除することを確認した。
